\chapter{Predictive versus descriptive information}\label{ch:predictive}

\Cref{ch:separation} ended with an empirical zero: on the six-label census, radius-two histograms leave no uncertainty about whether $\beta_2>0$. A zero residual entropy is a statement about a \emph{finite sample of signatures}. It can arise because the signature is informative, or because it is nearly injective, so that every system sits alone in its class. This chapter separates \emph{descriptive} information (how well the classes of a signature split the positives from the negatives in a sample) from \emph{predictive} information (how much a rule learned on one sample tells about topology in another). It also preserves the predictions that failed: a pre-registered claim that was falsified, and a general principle that was falsified by a certified counterexample. These are results, not embarrassments, and the later chapters depend on knowing which statements survived.

\section{Conventions}\label{sec:conv-info}
Let $\mathcal S$ be a finite set of systems (regimes), $Y:\mathcal S\to\{0,1\}$ the target $\mathbf 1[\beta_2>0]$ and $O$ an observation. A \emph{class} is a fibre of $O$; it is \emph{mixed} if it contains both values of $Y$. The \emph{plug-in} conditional entropy is computed from the empirical class frequencies of the sample,
\[
H_{\mathrm{plug}}(Y\mid O)=\sum_{c}\frac{|c|}{|\mathcal S|}\,\mathrm h\!\left(\frac{|c\cap Y^{-1}(1)|}{|c|}\right),\qquad \mathrm h(p)=-p\log_2p-(1-p)\log_2(1-p).
\]
Two ensembles are used. The \emph{natural} ensemble is uniform on regimes (positives are rare: $H(Y)=0.005713$ bits on the census). The \emph{balanced} ensemble weights positives and negatives equally, so $H(Y)=1$ bit and a rare positive is not drowned out. Both are experimental conventions, not probability laws of the grammar.
\begin{lemma}[a zero is cheap when signatures are fine]\label{lem:zero-cheap}
$H_{\mathrm{plug}}(Y\mid O)=0$ iff $O$ has no mixed class on $\mathcal S$. In particular, if $O$ is injective on $\mathcal S$ then $H_{\mathrm{plug}}(Y\mid O)=0$ for \emph{every} target $Y$. A zero residual therefore certifies information about $Y$ only relative to a null model that has the same class sizes.
\end{lemma}
\begin{proof}
A class contributes $\mathrm h(p)$, which vanishes iff $p\in\{0,1\}$ iff the class is pure. If $O$ is injective every class is a singleton.
\end{proof}
The source programme controlled for this with a \emph{random-placement null}: assign the positives at random among regimes of equal $N$ and count how many positives land in mixed classes of the signature under test (\cref{sec:separators}).

\section{Minimal separating signatures (descriptive)}\label{sec:separators}
Work on the census $\mathcal D$ with $15$ scalar features (counts, degree and ball-size moments, square counts) and the histograms $h_1,h_q,h_2$.
\begin{observation}[smallest separating signatures on the census]\label{obs:separators}
(a) No single scalar from the pool has zero residual; the best, $S_{BB}$, hides $36$ positives in $18$ mixed classes.
(b) The minimal separating \emph{pairs} include $(N,S_{BB})$, $(Q,S_{BB})$, $(S_B,S_{BB})$ and, in the second-order pool, $(Q,S_{s2})$, $(S_{d2},S_{s2})$, $(S_{d3},S_{s2})$, $(S_{q2},S_{s2})$, $(S_{s2},S_{ds})$, $(S_{s2},S_{dq})$, $(S_{s2},S_{sq})$. Every separating pair contains a second-order radius-two sum ($S_{BB}$ or $S_{s2}$). The full histogram $h_2$ and the $f$-vector also give zero residual.
(c) Of the $66$ positives left in mixed classes by $(N,h_1)$, adding \emph{one} scalar resolves: $S_{s2}$: all $66$; $S_{BB}$: all $66$; $f_3$: $64$; $S_{ds}$ and $S_{sq}$: $59$; $h_x$: $47$; $S_{q2},S_{dq}$: $44$; $Q$: $42$; $S_B$: $42$; nothing of first order ($E,S_{d2},S_{d3},d_{\max},d_{\min}$) resolves any.
(d) Non-triviality: for $(N,S_{BB})$ and the other separating pairs, $104$ of the $151$ positive masks lie in pure classes of size $\ge2$ (the largest pure positive class has $15$ members); placing the positives at random among regimes of equal $N$ hides $133$--$149$ of the $151$ (mean $143$) in mixed classes, and none of $300$ shuffles gave $0$. Forty-seven positives are alone in their class.
\status{\FiniteExact\ on the census (inherited, \texttt{NOT\_RERUN}); descriptive, no claim beyond the census}{GIG-INFO v003, R1, R2, R4}
\end{observation}

\section{A pre-registered test, and what it falsified}\label{sec:prereg}
The signatures $S_A=(N,S_{BB})$, $S_B=(Q,S_{s2})$ and $S_C=(S_{d2},S_{s2})$ were derived on the census as having zero mixed classes. Before any held-out system was generated, the following was registered (random seeds, sampling frame and feature definitions fixed, with the census features file hashed).
\begin{itemize}[nosep]
\item \textbf{H6} (held-out, six labels): $6000$ contents sampled uniformly from the six-subsets of $\{1,\dots,22\}$ whose maximum label lies in $17$--$22$ ($66{,}605$ such subsets), all regimes of each content.
\item \textbf{P1}: each of $S_A,S_B,S_C$ has zero mixed classes in H6 and in the pooled set $\mathcal D\cup$H6. \emph{Falsified by any mixed class.}
\item \textbf{P2}: the weaker signature $(N,h_1)$ has at least one positive in a mixed class in $\mathcal D\cup$H6.
\item \textbf{H7} (exploratory, seven labels): $120$ contents sampled from the seven-subsets of $\{1,\dots,14\}$; \emph{no prediction of success is made}.
\end{itemize}
\begin{observation}[outcome of the pre-registered held-out test]\label{obs:preregistered}
H6 contains $263{,}124$ regimes, $132$ positives; pooled with $\mathcal D$: $609{,}983$ regimes, $290$ positives.
\begin{itemize}[nosep]
\item $S_B=(Q,S_{s2})$: $0$ mixed classes in H6 and pooled. \textbf{Confirmed.}
\item $S_C=(S_{d2},S_{s2})$: $0$ mixed classes in H6 and pooled. \textbf{Confirmed.}
\item $(N,h_1,h_2)$ (not a registered prediction): $0$ mixed classes in H6 and pooled.
\item $S_A=(N,S_{BB})$: $3$ mixed classes ($3$ positives) in H6 and $6$ mixed classes ($10$ positives) pooled. \textbf{The pre-registered prediction for $S_A$ is falsified.} In $5$ of the $6$ classes the colliding systems differ in $Q$ and $S_{s2}$; in the sixth ($N=576$) they differ in $S_{s2}$ only.
\item P2: $(N,h_1)$ leaves $143$ of the $290$ pooled positives in mixed classes. \textbf{Confirmed.}
\end{itemize}
Consequently the registered claim ``all three pairs separate'' is false as stated (P1 fails through $S_A$). The weaker, more structured statement ``$(Q,S_{s2})$, $(S_{d2},S_{s2})$ and the $h_2$ histogram separate'' survived out of sample on $263$ thousand new regimes. Non-triviality on the pooled set: $206$ of the $290$ positives lie in pure classes of size $\ge2$; the random-placement null hides $267$--$285$ of $290$; the real value is $0$. In the seven-label exploratory sample ($10{,}130$ regimes, $11$ positives) every candidate signature, including $(N,h_1)$, has zero mixed classes while $N$ alone does not: the sample is too small and its signatures nearly injective to discriminate, and no transfer claim is made at $k=7$.
\status{\Bounded\ (inherited, \texttt{NOT\_RERUN}; pre-registration file hashed in the source manifest)}{GIG-INFO v003, R5, R6; PREREGISTRATION}
\end{observation}
$(N,S_{BB})$ is therefore a good fingerprint on $\mathcal D$ but not a law. This is the first of the two failed predictions this chapter preserves.

\section{A general principle that is false}\label{sec:principle-false}
After the pre-registered test the question became whether $(Q,S_{s2})$ separates for a structural reason --- for instance because all arithmetic systems are prefix-defined --- or only because of arithmetic closure.
\begin{observation}[the pair $(Q,S_{s2})$ is not sufficient in general]\label{obs:S2-falsified}
\begin{itemize}[nosep]
\item On arithmetic regimes ($\mathcal D\cup$H6, $609{,}983$ regimes, $290$ positives) there is no mixed class of $(Q,S_{s2})$, $(N,h_1,Q,S_{s2})$ or $(N,h_1,h_2)$, and every one of the $206$ positives that shares its $(Q,S_{s2})$ with another regime shares it only with positives (partner-positive rate $1.000$ against a base rate $0.0005$).
\item On \emph{prefix-family systems} (sets of orderings all of whose prefix sets lie in a family; $109{,}906$ distinct random systems, $4{,}796$ positives): $(Q,S_{s2})$ has $60$ mixed classes ($89$ positives hidden); $(N,Q,S_{s2})$ has $28$ ($52$ hidden); $(N,h_1,Q,S_{s2})$ has $24$ ($48$ hidden); $(N,h_1,Q,S_{s2},S_{ds})$ has $21$ ($43$).
\item \emph{A certified counterexample.} Two prefix-family systems with $N=388$, identical $h_1$, identical square-incidence histogram, $Q=397$, $S_{s2}=85536$, equal $S_{ds}$, $S_B$ and even $S_{BB}=96888$, but $\beta_2=1$ versus $0$ (the positive with an explicit $2$-cycle and cocycle verified from the definition; the negative with $\beta_2=0$ by two rank routes and an independent engine); they differ in $h_2$ and $f_3$ ($105$ vs $106$). Neither is an arithmetic census regime.
\item The full radius-two histogram $(N,h_1,h_2)$ has zero mixed classes in all three tested classes (arithmetic $609{,}983$; prefix families $109{,}906$ with $4{,}796$ positives; arbitrary vertex sets $30{,}000$ with $11{,}373$ positives); in prefix families $2{,}991$ positives share their $(N,h_1,h_2)$ with other systems, all partners positive.
\end{itemize}
Hence the separation by $(Q,S_{s2})$ is \emph{not} a consequence of the prefix structure or of local graph combinatorics alone, and no radius-two \emph{moment} signature $(N,Q,S_{s2},S_{ds},S_B,S_{BB})$ is sufficient in general. It holds on every arithmetic regime tested. Whether arithmetic closure forces it is open (\cref{prob:arith-origin}). The principle ``radius-two scalar moments determine topology'' is \textbf{falsified} as a general statement.
\status{\FiniteExact\ for the counterexample (certificate in the source release); \Bounded\ for the censuses (inherited, \texttt{NOT\_RERUN})}{GIG-INFO v004, R-A1}
\end{observation}
\begin{observation}[the seven hard positives]\label{obs:hard-positives}
On the census the positives left in mixed classes by $(N,h_1,S_{ds})$ number $7$ (by $(Q,S_{ds})$: $10$; by $(N,h_1,Q)$: $24$); on $\mathcal D\cup$H6 these numbers are $38,53,74$. In every one, $S_{s2}$ resolves the class, but narrowly: the positive's $S_{s2}$ is smaller than the partner's by $16$ (or $48$), relative margins of order $10^{-4}$ --- for instance content $(2,3,8,9,11,14)$ at $77$ and content $(2,3,4,8,14,15)$ at $1680$. In all $13$ mixed classes of $(N,h_1,Q,S_{ds})$ on $\mathcal D\cup$H6 the positive has the smaller $S_{s2}$ (gap $-16$ typically); in the $30$ mixed classes of $(N,h_1,Q)$ only $21$ have all positives below all negatives, so there is no one-sided inequality law. Why $S_{s2}$ --- a sum over pairs of two-step walks --- correlates with the topology is not explained; the simplest forbidden-subset motifs (dead subsets, complementary dead pairs) do not discriminate ($142/158$ positives against $1333/2000$ sampled negatives have a complementary dead pair).
\status{\FiniteExact\ (the margins) and \Bounded\ (the sign pattern); inherited}{GIG-INFO v004, R-A2, R-A3; v005, R-A4}
\end{observation}

\section{Descriptive versus predictive}\label{sec:oos}
\begin{observation}[plug-in information is descriptive]\label{obs:plugin}
On $\mathcal D\cup$H6, balanced ensemble ($609{,}983$ systems, $290$ positives): $H(Y\mid N)=0.4552$, $H(Y\mid N,h_1)=0.0126$, $H(Y\mid N,h_1,h_q)=0.0016$ and $H(Y\mid Q,S_{s2})=H(Y\mid N,h_1,Q,S_{s2})=H(Y\mid N,h_1,h_2)=0$ bits. On the natural ensemble over $\mathcal D$: $I(Y;N)=0.00124$ of $0.00571$ bits ($21.8\%$), $I(Y;N,h_1)=0.00509$ ($89.1\%$), zero residual for $(N,h_1,Q,S_{s2})$. These plug-in values measure how little the exact signature classes overlap, and fine signatures ($55{,}000$--$110{,}000$ classes) leave most negatives alone in their class (\cref{lem:zero-cheap}). They are \emph{not} predictive information.
\status{\FiniteExact\ with the stated warning (inherited, \texttt{NOT\_RERUN})}{GIG-INFO v004, R-B1}
\end{observation}
\begin{observation}[out-of-sample information]\label{obs:oos}
Define the information \emph{gained about $Y$} by an observation as $1-$ the balanced cross-entropy of a classifier fitted on one sample and evaluated on the other (the classifier family and the exact fitting protocol are those of the source release and are not reproduced here). Train on $\mathcal D$ and test on H6 (first number) or the reverse (second number), balanced test ensemble:
\begin{center}\small
\begin{tabular}{@{}lcc@{}}\toprule
observation & $\mathcal D\to$H6 & H6$\to\mathcal D$\\ \midrule
$N$ alone & $0.31$ & $0.38$\\
radius-one scalars & $0.28$ & $0.52$\\
radius-one $+$ square count & $0.44$ & $0.66$\\
all radius-two scalars & $0.44$ & $0.68$ (AUC $0.96$/$0.98$)\\
$(Q,S_{s2})$ alone & $0.29$ & $0.45$\\
exact-signature lookup with back-off & $-1.05$ & $-1.19$\\ \bottomrule
\end{tabular}
\end{center}
The exact lookup is worse than ignorance: $61\%$ of the held-out positives ($\mathcal D\to$H6) have a signature never seen in training. An observer learning from one sample therefore recovers at most about $0.4$--$0.7$ of the one bit about generative topology; the zero residual of exact signatures is a statement about collisions, not about generalisation.
\status{\Bounded\ (inherited, \texttt{NOT\_RERUN}); values are registered figures}{GIG-INFO v004, R-B2}
\end{observation}
\begin{figure}[ht]\centering\includegraphics[width=\textwidth]{figures/fig6_information_recovery.pdf}
\caption{(a) In-sample residual entropy $H(Y|O)$ of the hole indicator by observation level, balanced ensemble on $609{,}983$ six-label systems (log-style axis). (b) Information gained about $Y$ when a rule is trained on one sample and tested on the other. The two panels tell different stories: the in-sample residual falls to $0$ while the out-of-sample gain saturates near $0.4$--$0.7$ bits. \emph{Both panels plot registered figures from the source releases; none is recomputed here.}}\label{fig:recovery}
\end{figure}

\section{The same hierarchy for \texorpdfstring{$\beta_1$}{beta-1}}
\begin{observation}[like-for-like replication for the first Betti number]\label{obs:beta1-hierarchy}
With target $Y_1=\mathbf 1[\beta_1>0]$ the ordering $N\to$ radius one $\to$ radius two reproduces in both domains. Symbolic systems ($508{,}680$ systems, $11{,}547$ positives): $N$ leaves $123$ mixed classes ($10{,}009$ positives hidden), $(N,h_1)$ leaves $5$ ($34$ hidden), $(N,h_1,Q,S_{s2})$ and $(N,h_1,h_2)$ leave $0$. Arithmetic census ($346{,}859$ regimes, $2{,}807$ positives): $N$: $262$ mixed; $(N,h_1)$: $120$ mixed ($233$ hidden); $(N,h_1,Q)$: $13$; $(N,h_1,Q,S_{s2})$ and $(N,h_1,h_2)$: $0$. Radius-one information is nearly sufficient for $\beta_1$ in symbolic systems but not in arithmetic ones; the square count and $S_{s2}$ (or $h_2$) close the gap in both. Plug-in, in-sample, with the caveat of \cref{obs:plugin}.
\status{\FiniteExact\ and \Bounded\ (inherited, \texttt{NOT\_RERUN})}{GIG-INFO v004, R-C3}
\end{observation}
Observations~\ref{obs:preregistered} and \ref{obs:S2-falsified} are why the thesis of this manuscript is stated in terms of \emph{structure} (forbidden histories, locality), not in terms of a statistic: the statistics that work do so on the families where they were found, and the mechanism that makes them work is not known. The next part supplies the structural mechanism that \emph{is} proved.
